The code used for this work can be found in the following GitHub repository: \url{https://github.com/TheophileLaPelouse/projet_g3/tree/main/python}. The code is organized as a python module named \texttt{commu\_opti}. It is structured as the table \ref{tab:code_structure} shows. Then a \texttt{test.py} file is used to create the first tests presented and a file \texttt{result.py} is used to create the results presented.


\begin{table}[h]
    \centering
    \begin{tabular}{|c|c|c|c|}
        \hline
        \multicolumn{3}{|c|}{} & \textbf{Description} \\
        % \hline
        % & &  & \textbf{Explanation} \\
        \hline
        \textbf{commu\_opti} & & & Global python module \\
        \cline{2-4}
        & \texttt{commu\_builder.py} & & Functions for building communities \\
        \cline{2-4}
        & \textbf{community} & &  \\

        & & \texttt{device.py} & Define device classes and models \\
        \cline{3-4}
        & & \texttt{member.py} & Define member class and model \\
        \cline{3-4}
        & & \texttt{community.py} & Define community class and model \\
        \cline{3-4}
        & & \texttt{utils.py} & None class specific functions \\
        \cline{3-4}
        \cline{2-4}
        & \textbf{opti} & &  \\

        & & \texttt{solving.py} & Solve Pyomo optimization models \\
        \cline{3-4}
        \cline{2-4}
        & \textbf{data} & &  \\

        & & \texttt{generate\_data.py} & Generate synthetic data \\
        \cline{3-4}
        \cline{2-4}
        & \textbf{plotting} & &  \\

        & & \texttt{plot\_functions.py} & Define plotting functions \\
        \hline
    \end{tabular}
    \caption{Structure of the code developed}
    \label{tab:code_structure}
\end{table}

